% =====================================================
% CONCLUSION
% =====================================================
\chapter*{Conclusion}
\addcontentsline{toc}{chapter}{Conclusion}

This project implemented a complete supervised Machine Learning pipeline applied to network traffic anomaly detection: exploratory analysis of a dataset containing 103,000 observations, preprocessing (cleaning, imputation, encoding), training of a \textbf{Random Forest} model (100 trees), evaluation, and validation.

\section*{Results}
\addcontentsline{toc}{section}{Results}

On the test set (30,900 samples): \textbf{99.99\% accuracy}, precision/recall/F1-score = 1.00 for both classes, with only 2 false positives and 0 false negatives. Feature importance analysis identified the most discriminative network indicators.

\section*{Limitations}
\addcontentsline{toc}{section}{Limitations}

\begin{itemize}[leftmargin=*,itemsep=2pt]
    \item A near-perfect score should be interpreted with caution: it may reflect a very clear separation specific to this dataset and may not generalize to real-world traffic.
    \item Validation on a second independent dataset was planned (Part 4) but not executed due to the lack of an additional dataset.
    \item GNS3 integration (Part 5) could not be completed due to hardware constraints.
    \item Mean imputation is a simplification that may smooth out important signals.
\end{itemize}

\section*{Future work}
\addcontentsline{toc}{section}{Future work}

\begin{itemize}[leftmargin=*,itemsep=2pt]
    \item Test the model on larger and more diverse datasets.
    \item Compare with other algorithms (XGBoost, SVM, logistic regression).
    \item Explore deep learning approaches (LSTM, autoencoders) for sequential traffic analysis.
    \item Integrate the model into a real-time processing pipeline (live traffic monitoring).
    \item Apply cross-validation and hyperparameter tuning.
    \item Complete GNS3 integration on a suitable platform.
\end{itemize}

